\section{Discussion}
We next discuss concurrent works related to ours and highlight their key differences. Rolling Forcing~\cite{liu2025rollingforcingautoregressivelong} extends the concept of Rolling Diffusion~\cite{ruhe2024rollingdiffusionmodels} by applying progressively varied noise levels to different video frames. It integrates attention sink frames for balancing short- and long-term consistency, while training efficiency is improved by sampling non-overlapping frames. LongLive~\cite{yang2025longlive} builds upon Self Forcing~\cite{huang2025self}, introducing KV re-caching for prompt switching and leveraging clean contexts. It further employs attention sink frames to mitigate error accumulation by repeatedly applying DMD to future frames beyond the teacher’s horizon. Our approach is most closely related to LongLive where we also incorporate DMD into long self-rolled sequences in a windowed fashion with clean context as detailed in~\cref{sec.extended_dmd}. Unlike LongLive, however, our simplified design avoids reliance on attention sink frames to counter error accumulation, which was shown to be a key design of LongLive.

\textbf{Both Rolling Forcing~\cite{liu2025rollingforcingautoregressivelong} and LongLive~\cite{yang2025longlive}, as well as our method}, are able to generate high-quality videos up to several minutes long, which marks a significant advance in autoregressive long video generation compared to previous methods.

